\documentclass{article}
\usepackage[T1]{fontenc}
\usepackage{lmodern}
\usepackage{graphicx}
\usepackage{amsmath,amssymb}
\usepackage[a4paper,margin=2cm]{geometry}
\usepackage[absolute,overlay]{textpos}
\setlength{\TPHorizModule}{1mm}
\setlength{\TPVertModule}{1mm}
\usepackage[indonesian]{babel}
\usepackage{hyperref}
\setlength{\emergencystretch}{2em}
% unit: Halaman 5

\begin{document}

% === Halaman 5 (python/native) ===

Akar Primitif dan Logaritma Diskrit

• Jika n adalah bilangan bulat, maka a disebut akar primitif dari n jika perpangkatan 

a, a2, …, a(n)

(dalam sebuah modulus n) 

    menghasilkan nilai yang berbeda dan semuanya relatif prima dengan n.

• Secara khusus, jika p adalah bilangan prima, maka a disebut akar primitif dari p 

jika perpangkatan 

a, a2, …, ap – 1

(dalam modulus p)

menghasilkan nilai-nilai yang berbeda 

(ingatlah dari fungsi toitient Euler, bahwa jika p prima maka (p) = p – 1) 

5

\end{document}
